\documentclass[12pt]{exam}
\usepackage{amsthm}
\usepackage{libertine}
\usepackage[utf8]{inputenc}
\usepackage[margin=1in]{geometry}
\usepackage{amsmath,amssymb}
\usepackage{multicol}
\usepackage[shortlabels]{enumitem}
\usepackage{siunitx}
\usepackage{cancel}
\usepackage{graphicx}
\usepackage{pgfplots}
\usepackage{listings}
\usepackage{tikz}


\pgfplotsset{width=10cm,compat=1.9}
\usepgfplotslibrary{external}
\tikzexternalize

\newcommand{\class}{APCS} % This is the name of the course 
\newcommand{\examnum}{HW1} % This is the name of the assignment
\newcommand{\examdate}{9-2-26} % This is the due date
\newcommand{\timelimit}{}





\begin{document}
\pagestyle{plain}
\thispagestyle{empty}

\noindent
\begin{tabular*}{\textwidth}{l @{\extracolsep{\fill}} r @{\extracolsep{6pt}} l}
\textbf{\class} & 
\textbf{Name:} & \textit{Toby Semple}\\ %Your name here instead, obviously 
\textbf{\examnum} &&\\
\textbf{\examdate} &&\\
\end{tabular*}\\
\rule[2ex]{\textwidth}{2pt}
% ---




\begin{enumerate} %You can make lists!

\item 32,767 = 31 * 1,057 \begin{enumerate}
    \item 2. It is both a number greater than one and is able to factor into the given number.

\end{enumerate}

\item \begin{enumerate}
    \item There are no values of 'n' for which $3^n - 1$ is a prime.
    \item 
\begin{table}[h]
    \centering % Centers the table on the page
    \caption{An example table using booktabs.}
    \label{tab:example_table}
    
    \begin{tabular}{lccr} % 3 columns: left-aligned, centered, right-aligned
        \toprule
        \textbf{x} & \textbf{Is it prime?} & \textbf{$3^x - 1$} & \textbf{Is it Prime?}\\
        \midrule
        1            & yes               & \$2        & no (even) \\
        2            & yes               & \$8        & no (even)\\
        3            & yes               & \$26       & no (even)\\
        4            & no                & \$80       & no (even)\\
        5            & yes               & \$242      & no (even)\\
        6            & no                & \$728      & no (even)\\
        7            & yes               & \$2186     & no (even)\\
        8            & no                & \$6560     & no (even)\\
        9            & yes               & \$19682    & no (even)\\
        10           & no                & \$59048    & no (even)\\

        \bottomrule
    \end{tabular}
\end{table}



    \item When n is prime, then $3^n - 2^n$ might be prime.

    \begin{table}[h]
    \centering % Centers the table on the page
    \caption{An example table using booktabs.}
    \label{tab:example_table}
    
    \begin{tabular}{lccr} % 3 columns: left-aligned, centered, right-aligned
        \toprule
        \textbf{x} & \textbf{Is it prime?} & \textbf{$3^x - 2^n$} & \textbf{Is it Prime?}\\
        \midrule
        1            & yes               & \$1        & yes                  \\
        2            & yes               & \$5        & yes                  \\
        3            & yes               & \$19       & yes                  \\
        4            & no                & \$65       & no (divisible by 5)  \\
        5            & yes               & \$211      & yes                  \\
        6            & no                & \$665      & no (divisible by 5)  \\
        7            & yes               & \$2059     & no (divisible by 29) \\
        8            & no                & \$6305     & no (divisible by 5)  \\
        9            & yes               & \$19171    & no (divisible by 19) \\
        10           & no                & \$58025    & no (divisible by 5)  \\

        \bottomrule
    \end{tabular}
\end{table}
\end{enumerate}


\item \begin{enumerate}
    \item $(2*3*5*7) + 1 = 211$
    This is a prime number because there are no factors below $sqrt(211) < 15$ except 1.
    \item $(2*5*11) + 1 = 111$ // $111 / 3 = 37$
    37 is a prime number because there are no factors below $sqrt(37) < 7$ except 1.
\end{enumerate}

\item $(5+1)! + 2 = 720$ 
The requested sequence is: 722, 723, 724, 725, & 726
$722 / 2 = 361
723 / 3 = 241
724 / 4 = 181
725 / 5 = 145
726 / 6 = 121$

\item No, because that would require one of the numbers to be a multiple of three. 
If the first number 'x' would be such that $x / 3 = y r 1$, then the next number must be a multiple of three, $( 1 + 2 = 3,)$ gets rid of the remainder.
If the first number is such that the remainder would be two, the next one would have a remainder of 1 $(2 + 2 = 4, 4 / 3 = 1 r 1)$, the third one would then be a multiple of 3.

\item 



\end{enumerate}

Look! Math! $f(x)=4x^2 - \frac{3}{4}x +1$ %use the dollar signs to go into math mode


\end{document}